\section{Related Work}
\label{sec:related}

This section situates \PAR among prior work across three axes: formal agent models,
agent frameworks, and type-safe runtime systems. We highlight the gaps each prior
approach leaves unfilled and that \PAR aims to close.

\subsection{Formal Agent Models}
\label{ssec:formal-models}

The most directly relevant formal work is $\lambda_A$~\cite{lambda-a}, a Coq
formalization establishing 42 theorems on the typed composition of agents. $\lambda_A$
proves that composite agents preserve safety invariants under parallel composition.
However, the formalization is not runnable: it serves as a specification, not an
implementation. \PAR differs along several dimensions. First, we provide a concrete
executable implementation in OCaml, not a proof-only artifact. Second, we introduce a
tool \ADT\ with sealed command-response contracts, going beyond the untyped tool
invocations in $\lambda_A$. Third, our middleware model imposes a well-defined
request-response lifecycle that $\lambda_A$ leaves implicit. Fourth, \PAR tracks
resource bounds (memory, tool calls, step limits) as part of the state machine;
$\lambda_A$ does not bound resources. Despite these differences, $\lambda_A$ remains
the closest related formal treatment of agent composition typing and we build on its
typed composition theorems.

CTEGs~\cite{ctegs} model recursive agent execution traces as causal-temporal graphs,
capturing dependencies between agent actions and their temporal ordering. This work is
complementary to ours: CTEGs focus on trace semantics and causal reasoning, whereas
\PAR focuses on the runtime lifecycle and state machine enforcement during execution.
We do not model trace causality, but \PAR's observation logs could in principle be
used to reconstruct execution traces.

The Model Context Protocol has a formal treatment via process calculus
~\cite{mcp-formal}, which models MCP messages as typed processes with session
semantics. This work is orthogonal to ours: it specifies the protocol layer, not the
runtime state machine that interprets protocol messages at an agent's call site. \PAR
could consume MCP-compliant tools; we do not depend on a specific protocol version.

\subsection{Agent Frameworks}
\label{ssec:frameworks}

Several projects provide practical agent infrastructure without formal guarantees.
The \texttt{oharness\_tools} library (Rust) introduces a trait-based tool \ADT\
allowing developers to register commands with typed inputs and outputs. However, it
provides no state machine abstraction, no DSL for agent policies, and no formal
semantics beyond trait bounds. The system is runnable but does not enforce transition
safety between agent states.

The \texttt{tower-llm} framework (Rust) offers a Service-trait middleware model for
LLM request/response processing, enabling composable pipeline stages such as
retrieval augmentation or response validation. This middleware concept aligns with
\PAR's request middleware, but \texttt{tower-llm} lacks a formal composition proof
and does not verify that middleware chains preserve agent invariants. Tool calls in
\tower-llm\ are untyped Rust functions, not a sealed ADT with protocol-level
contract guarantees.

The \texttt{llaml} project (OCaml) provides partial tool typing via OCaml functors,
giving some structural safety for tool registration. It does not, however, provide a
state machine, a DSL for policy specification, or a middleware model. Its
executability is limited to single-shot tool calls without lifecycle management.

Microsoft's Semantic Kernel~\cite{semantic-kernel} (C\#) implements a middleware
pipeline via dependency injection, allowing planners and memory plugins to be wired
into a kernel. The pipeline is extensible but carries no formal semantics. Like
\tower-llm, middleware correctness is trusted to the developer; \PAR proves
middleware preserves invariant transitions.

\subsection{Type-Safe Runtime Systems}
\label{ssec:type-safe-runtimes}

Algebraic effects in OCaml 5~\cite{algebraic-effects-ocaml} provide a foundation for
structured concurrency with typed effect signatures. Our work currently uses ADTs
and pattern matching for state transitions rather than effect handlers, though we see
potential in encoding \PAR's middleware as effect handlers in future work. The
current ADT-based approach is sufficient for our formal guarantees and avoids the
complexity of effect inference.

Monadic Context Engineering~\cite{monadic-context} proposes context management via
monads, encoding request-scoped values as monadic carry-along state. This is
complementary to our approach: we use the \Context record for request-scoped values
rather than a monad, trading monad laws for a flat, inspectable record that can be
logged and introspected at runtime.

Effect Systems for AI Runtimes~\cite{effect-ai} apply effect typing to agent
runtimes, tracking which tools a given agent step may invoke. The scope differs from
ours: effect systems enforce that calls stay within a permitted effect set, while
\PAR enforces that transitions stay within a permitted state graph. Both aim for
type-safe agency, but the guarantees are orthogonal: effect safety versus state
machine safety.

\begin{table}[t]
\caption{Comparison of agent runtimes and formal models.}
\label{tab:comparison}
\centering
\small
\begin{tabular}{l c c c c c c}
\toprule
Project & Lang. & Tool ADT & State Mach. & DSL & Middleware & Runnable \\
\midrule
$\lambda_A$~\cite{lambda-a} & Coq & \ding{55} & \ding{51} & \ding{51} & \ding{55} & \ding{55} \\
CTEGs~\cite{ctegs} & --- & \ding{55} & \ding{55} & \ding{51} & \ding{55} & \ding{55} \\
oharness\_tools & Rust & \ding{51} & \ding{55} & \ding{55} & Observer & \ding{51} \\
tower-llm & Rust & \ding{55} & \ding{55} & \ding{55} & \ding{51} & \ding{51} \\
llaml & OCaml & partial & \ding{55} & \ding{55} & \ding{55} & \ding{51} \\
Semantic Kernel~\cite{semantic-kernel} & C\# & \ding{55} & \ding{55} & \ding{55} & \ding{51} & \ding{51} \\
\textbf{P-A-R (ours)} & \textbf{OCaml} & \ding{51} & \ding{51} & \ding{51} & \ding{51} & \ding{51} \\
\bottomrule
\end{tabular}
\end{table}

\PAR uniquely occupies the intersection of formal specification and runnable
implementation. Table~\ref{tab:comparison} summarizes the tradeoffs: most prior work
excels in either expressiveness (formal models) or executability (frameworks), but
not both. The few systems that are both runnable and typed lack a state machine, a
DSL, or middleware. \PAR bridges this gap by providing a Coq-adjacent formal
framework (typed tool ADT, middleware invariants, transition safety) inside a
concrete OCaml runtime that can be executed, tested, and deployed today.